\documentclass[11pt,a4paper]{article}
\usepackage{shared/luxcover}
\usepackage[utf8]{inputenc}
\usepackage[T1]{fontenc}
\usepackage{amsmath,amssymb,amsfonts,amsthm}
\usepackage{graphicx}
\usepackage{hyperref}
\usepackage{booktabs}
\usepackage{listings}
\input{shared/lstlang}
\usepackage{xcolor}
\usepackage{enumitem}
\usepackage{float}
\usepackage{tikz}
\usetikzlibrary{arrows.meta,positioning}

\newtheorem{theorem}{Theorem}
\newtheorem{lemma}[theorem]{Lemma}
\newtheorem{proposition}[theorem]{Proposition}
\newtheorem{definition}{Definition}
\newtheorem{corollary}[theorem]{Corollary}
\newtheorem{remark}{Remark}
\newtheorem{example}{Example}

\lstset{
    basicstyle=\small\ttfamily,
    keywordstyle=\color{blue}\bfseries,
    commentstyle=\color{gray},
    stringstyle=\color{red},
    breaklines=true,
    frame=single,
    numbers=left,
    numberstyle=\tiny\color{gray}
}

\hypersetup{
    colorlinks=true,
    linkcolor=blue,
    citecolor=blue,
    urlcolor=blue
}

\title{Lattice Cryptography from First Principles:\\A Builder's Guide}
\author{
Lux Industries, Inc.\thanks{Corresponding author: research@lux.network}\\
\texttt{research@lux.network}
}
\date{2025}

\begin{document}

\luxcoverpage

\begin{abstract}
Lattice-based cryptography is the foundation of post-quantum security: ML-DSA (Dilithium)
for signatures, ML-KEM (Kyber) for key encapsulation, and Ringtail for threshold schemes
all build on the hardness of lattice problems. This paper teaches lattice cryptography
from the ground up---not as a survey of results, but as a builder's guide for engineers
implementing these schemes. We cover: lattices and their geometry, the Learning With Errors
(LWE) problem and why it is hard, Ring-LWE and Module-LWE for practical efficiency, the
Number Theoretic Transform (NTT) for polynomial arithmetic, how ML-DSA and ML-KEM work
step by step, and how Ringtail extends these to threshold operations. We include Go code
for every non-trivial algorithm and explain every design decision in terms of the
security/performance tradeoff it represents. This is the paper we wish existed when we
started implementing post-quantum cryptography for Lux Network.
\end{abstract}

\tableofcontents
\newpage

%% ============================================================
\section{What Is a Lattice}

\begin{definition}[Lattice]
A lattice $\Lambda$ in $\mathbb{R}^n$ is the set of all integer linear combinations of
a set of linearly independent vectors $\mathbf{b}_1, \ldots, \mathbf{b}_n \in \mathbb{R}^n$:
\[
\Lambda = \left\{ \sum_{i=1}^n z_i \mathbf{b}_i \;:\; z_i \in \mathbb{Z} \right\}
\]
The vectors $\{\mathbf{b}_i\}$ form a \emph{basis} for $\Lambda$.
\end{definition}

Think of a lattice as a grid---a regular arrangement of points in space. In two dimensions,
it looks like graph paper tilted and stretched. The crucial property: the grid is
\emph{discrete} (points are separated by gaps) but \emph{infinite} (it extends in all
directions).

\subsection{Why Lattices for Cryptography}

Classical public-key cryptography (RSA, elliptic curves) relies on the hardness of
factoring or discrete logarithms. Shor's algorithm solves both in polynomial time on a
quantum computer. Lattice problems---finding short vectors in high-dimensional
lattices---resist all known quantum algorithms.

The two fundamental hard problems:

\begin{definition}[Shortest Vector Problem (SVP)]
Given a lattice basis $\mathbf{B}$, find the shortest nonzero lattice vector.
\end{definition}

\begin{definition}[Closest Vector Problem (CVP)]
Given a lattice basis $\mathbf{B}$ and a target point $\mathbf{t}$, find the lattice
vector closest to $\mathbf{t}$.
\end{definition}

Both are NP-hard in their exact forms, and the best known algorithms (classical or
quantum) for approximate versions run in exponential time in the lattice dimension $n$.

\subsection{The Geometry of Hard Problems}

In low dimensions ($n = 2, 3$), finding short vectors is easy---look at the basis and
reduce it. In high dimensions ($n = 256, 512, 1024$), the geometry becomes pathological:
the shortest vector is exponentially shorter than the basis vectors, but finding it
requires exploring an exponentially large search space. This is the core intuition:
lattice cryptography works because high-dimensional geometry is hard.

%% ============================================================
\section{Learning With Errors (LWE)}

\subsection{The Problem}

\begin{definition}[Learning With Errors]
Given $m$ samples of the form $(\mathbf{a}_i, b_i)$ where:
\begin{itemize}[nosep]
  \item $\mathbf{a}_i \in \mathbb{Z}_q^n$ is uniformly random,
  \item $b_i = \langle \mathbf{a}_i, \mathbf{s} \rangle + e_i \pmod{q}$,
  \item $\mathbf{s} \in \mathbb{Z}_q^n$ is a secret vector,
  \item $e_i$ is a small ``error'' sampled from a discrete Gaussian,
\end{itemize}
recover $\mathbf{s}$.
\end{definition}

Without the error $e_i$, this is a system of linear equations---solvable by Gaussian
elimination. The error makes the problem hard: it turns an easy linear algebra problem
into a problem equivalent to finding short vectors in a lattice.

\subsection{Why LWE Is Hard}

Regev (2005) proved that solving LWE is at least as hard as solving worst-case lattice
problems (approximate SVP and approximate SIVP) in dimension $n$. This is a
\emph{worst-case to average-case} reduction: even random LWE instances are as hard as
the hardest lattice instances. No other family of cryptographic assumptions has this
property.

\subsection{LWE as Encryption}

LWE directly yields a public-key encryption scheme:

\begin{itemize}[nosep]
  \item \textbf{Key generation}: Choose secret $\mathbf{s}$. Publish samples
    $(\mathbf{A}, \mathbf{b} = \mathbf{A}\mathbf{s} + \mathbf{e})$.
  \item \textbf{Encryption of bit $m$}: Choose random $\mathbf{r}$. Compute
    $\mathbf{u} = \mathbf{A}^T \mathbf{r}$ and $v = \mathbf{b}^T \mathbf{r} + m \cdot \lfloor q/2 \rfloor$.
  \item \textbf{Decryption}: Compute $v - \mathbf{u}^T \mathbf{s} = m \cdot \lfloor q/2 \rfloor + \text{small noise}$.
    Round to recover $m$.
\end{itemize}

The decryption works because the noise terms are small relative to $q/2$. The scheme is
secure because distinguishing $(\mathbf{A}, \mathbf{A}\mathbf{s} + \mathbf{e})$ from
uniform random is the LWE problem.

%% ============================================================
\section{Ring-LWE and Module-LWE}

\subsection{The Efficiency Problem}

Plain LWE has a fundamental efficiency issue: the public key is a matrix $\mathbf{A} \in
\mathbb{Z}_q^{m \times n}$, which has size $O(n^2)$. For 128-bit security, $n \approx 512$,
so keys are hundreds of kilobytes. This is impractical for blockchain use.

\subsection{Ring-LWE}

Ring-LWE replaces vectors over $\mathbb{Z}_q$ with elements of a polynomial ring:

\[
R_q = \mathbb{Z}_q[X] / (X^n + 1)
\]

where $n$ is a power of 2 and $q$ is prime. An element of $R_q$ is a polynomial of
degree $< n$ with coefficients in $\mathbb{Z}_q$.

\begin{definition}[Ring-LWE]
Given samples $(a_i, b_i = a_i \cdot s + e_i) \in R_q \times R_q$, recover $s \in R_q$.
\end{definition}

The key advantage: a single ring element replaces an entire vector, reducing key sizes
from $O(n^2)$ to $O(n)$. Polynomial multiplication in $R_q$ replaces matrix-vector
multiplication, and NTT makes it $O(n \log n)$ instead of $O(n^2)$.

\subsection{Module-LWE}

Module-LWE (M-LWE) is a middle ground: it uses \emph{small matrices of ring elements}
rather than single ring elements or large integer matrices.

\begin{definition}[Module-LWE]
Let $k$ be a small integer (2, 3, or 4). Given
$(\mathbf{A}, \mathbf{b} = \mathbf{A}\mathbf{s} + \mathbf{e})$ where
$\mathbf{A} \in R_q^{k \times k}$, $\mathbf{s}, \mathbf{e} \in R_q^k$,
recover $\mathbf{s}$.
\end{definition}

M-LWE is the basis for ML-KEM (Kyber) and ML-DSA (Dilithium). The module dimension $k$
parameterizes the security level:

\begin{table}[H]
\centering
\small
\begin{tabular}{@{}lcccc@{}}
\toprule
\textbf{Scheme} & $n$ & $k$ & $q$ & \textbf{Security} \\
\midrule
ML-KEM-512 (Kyber-512) & 256 & 2 & 3329 & 128-bit \\
ML-KEM-768 (Kyber-768) & 256 & 3 & 3329 & 192-bit \\
ML-KEM-1024 (Kyber-1024) & 256 & 4 & 3329 & 256-bit \\
ML-DSA-44 (Dilithium2) & 256 & 4/4 & 8380417 & 128-bit \\
ML-DSA-65 (Dilithium3) & 256 & 6/5 & 8380417 & 192-bit \\
ML-DSA-87 (Dilithium5) & 256 & 8/7 & 8380417 & 256-bit \\
\bottomrule
\end{tabular}
\caption{NIST post-quantum standard parameters. All use $n = 256$.}
\end{table}

%% ============================================================
\section{The Number Theoretic Transform}

Polynomial multiplication in $R_q$ is the bottleneck operation. Naive multiplication
is $O(n^2)$. The NTT reduces this to $O(n \log n)$.

\subsection{What Is NTT}

The NTT is the discrete Fourier transform over a finite field $\mathbb{Z}_q$ instead of
$\mathbb{C}$. It transforms a polynomial from coefficient representation to evaluation
representation (values at the $n$-th roots of unity in $\mathbb{Z}_q$).

\begin{definition}[NTT]
Given a polynomial $a(X) = \sum_{i=0}^{n-1} a_i X^i$ and a primitive $n$-th root of
unity $\omega \in \mathbb{Z}_q$:
\[
\hat{a}_j = \text{NTT}(a)_j = \sum_{i=0}^{n-1} a_i \omega^{ij} \pmod{q}
\]
\end{definition}

In evaluation representation, polynomial multiplication becomes pointwise:
\[
\text{NTT}(a \cdot b)_j = \text{NTT}(a)_j \cdot \text{NTT}(b)_j
\]

So: transform both polynomials ($2 \times O(n \log n)$), multiply pointwise ($O(n)$),
inverse-transform the result ($O(n \log n)$). Total: $O(n \log n)$.

\subsection{The Butterfly Operation}

The NTT is computed using the Cooley-Tukey butterfly, the same structure as the FFT:

\begin{lstlisting}[language=Go,caption={NTT butterfly operation in Go}]
// Butterfly: the core NTT operation
// In-place: a[j], a[j+step] are updated
func butterfly(a []uint32, j, step int, w uint32, q uint32) {
    t := mulmod(w, a[j+step], q)
    a[j+step] = submod(a[j], t, q)
    a[j] = addmod(a[j], t, q)
}

// Full NTT (iterative, in-place, Cooley-Tukey)
func NTT(a []uint32, omegas []uint32, q uint32) {
    n := len(a)
    bitReverse(a) // Reorder elements

    for step := 1; step < n; step <<= 1 {
        for j := 0; j < n; j += step << 1 {
            for k := 0; k < step; k++ {
                butterfly(a, j+k, step, omegas[step+k], q)
            }
        }
    }
}
\end{lstlisting}

\subsection{Montgomery Multiplication}

Modular multiplication ($a \cdot b \mod q$) is expensive because division is slow.
Montgomery multiplication avoids division by working in Montgomery form:

\begin{lstlisting}[language=Go,caption={Montgomery multiplication}]
const (
    Q     uint32 = 3329       // Kyber modulus
    QInv  uint32 = 62209      // Q^{-1} mod 2^16
    Mont  uint32 = 2285       // 2^16 mod Q (Montgomery constant)
)

// Montgomery reduction: given a < Q * 2^16,
// compute a * 2^{-16} mod Q without division
func montgomeryReduce(a uint32) uint16 {
    t := uint16(a) * uint16(QInv)
    return uint16((a - uint32(t)*uint32(Q)) >> 16)
}

// Multiply in Montgomery form
func mulMont(a, b uint16) uint16 {
    return montgomeryReduce(uint32(a) * uint32(b))
}
\end{lstlisting}

%% ============================================================
\section{How ML-KEM (Kyber) Works}

ML-KEM is a key encapsulation mechanism (KEM): it produces a shared secret between two
parties. It builds on M-LWE.

\subsection{Key Generation}

\begin{enumerate}[nosep]
  \item Sample random matrix $\mathbf{A} \in R_q^{k \times k}$ (from a seed via SHAKE-128).
  \item Sample secret $\mathbf{s} \in R_q^k$ and error $\mathbf{e} \in R_q^k$ with small
    coefficients (centered binomial distribution).
  \item Compute $\mathbf{t} = \mathbf{A}\mathbf{s} + \mathbf{e}$.
  \item Public key: $(\text{seed for } \mathbf{A}, \mathbf{t})$. Secret key: $\mathbf{s}$.
\end{enumerate}

\begin{lstlisting}[language=Go,caption={ML-KEM key generation (simplified)}]
func KeyGen() (PublicKey, SecretKey) {
    seed := randBytes(32)
    A := expandA(seed)  // k x k matrix of ring elements

    s := sampleCBD(eta1) // Secret: small coefficients
    e := sampleCBD(eta1) // Error: small coefficients

    // t = A*s + e (all operations in NTT domain)
    sNTT := nttVector(s)
    t := addVectors(mulMatVec(A, sNTT), nttVector(e))

    pk := PublicKey{Seed: seed, T: compress(t)}
    sk := SecretKey{S: sNTT}
    return pk, sk
}
\end{lstlisting}

\subsection{Encapsulation}

\begin{enumerate}[nosep]
  \item Generate random message $m \gets \{0,1\}^{256}$.
  \item Derive randomness: $(r, K) = G(m \| H(\text{pk}))$.
  \item Encrypt $m$ under the public key using randomness $r$:
    $\mathbf{u} = \mathbf{A}^T \mathbf{r} + \mathbf{e}_1$ and
    $v = \mathbf{t}^T \mathbf{r} + e_2 + \lceil q/2 \rceil \cdot m$.
  \item Ciphertext: $(\mathbf{u}, v)$. Shared secret: $K$.
\end{enumerate}

\subsection{Decapsulation}

\begin{enumerate}[nosep]
  \item Decrypt: $m' = \text{Decompress}(v - \mathbf{s}^T \mathbf{u})$.
  \item Re-encrypt: $(\mathbf{u}', v') = \text{Enc}(\text{pk}, m'; r')$.
  \item If $(\mathbf{u}', v') = (\mathbf{u}, v)$, output $K = \text{KDF}(m' \| H(\text{ct}))$.
  \item Otherwise, output a random key (implicit rejection).
\end{enumerate}

The re-encryption step (Fujisaki-Okamoto transform) converts CPA security to CCA
security: it prevents the adversary from learning anything by submitting crafted
ciphertexts.

%% ============================================================
\section{How ML-DSA (Dilithium) Works}

ML-DSA is a digital signature scheme based on the ``Fiat-Shamir with aborts'' paradigm
over Module-LWE.

\subsection{Key Generation}

\begin{enumerate}[nosep]
  \item Sample $\mathbf{A} \in R_q^{k \times \ell}$ from a seed.
  \item Sample secret vectors $\mathbf{s}_1 \in R_q^\ell$, $\mathbf{s}_2 \in R_q^k$
    with small coefficients.
  \item Compute $\mathbf{t} = \mathbf{A}\mathbf{s}_1 + \mathbf{s}_2$.
  \item Public key: $(\text{seed}, \mathbf{t})$. Secret key: $(\mathbf{s}_1, \mathbf{s}_2)$.
\end{enumerate}

\subsection{Signing (Fiat-Shamir with Aborts)}

\begin{enumerate}[nosep]
  \item Sample random masking vector $\mathbf{y}$ with bounded coefficients.
  \item Compute $\mathbf{w} = \mathbf{A}\mathbf{y}$ and extract high bits: $\mathbf{w}_1$.
  \item Compute challenge: $c = H(\text{msg} \| \mathbf{w}_1)$.
  \item Compute response: $\mathbf{z} = \mathbf{y} + c \cdot \mathbf{s}_1$.
  \item \textbf{Rejection sampling}: If $\|\mathbf{z}\|_\infty$ is too large, \emph{restart}.
    This prevents the signature from leaking information about $\mathbf{s}_1$.
  \item Output signature: $(c, \mathbf{z})$.
\end{enumerate}

The rejection sampling step is critical: without it, statistical analysis of many
signatures reveals $\mathbf{s}_1$. With it, the distribution of $\mathbf{z}$ is
independent of $\mathbf{s}_1$ (up to the bound check).

\begin{lstlisting}[language=Go,caption={ML-DSA signing with rejection sampling}]
func Sign(sk *SecretKey, msg []byte) Signature {
    for {
        // Sample masking vector
        y := sampleMask(gamma1)

        // w = A * y
        w := mulMatVec(sk.A, nttVector(y))
        w1 := highBits(w, alpha)

        // Challenge
        c := challengeHash(msg, w1)
        cPoly := sampleChallenge(c)

        // Response
        z := addVectors(y, scalarMulVector(cPoly, sk.S1))

        // Rejection: check norm bound
        if infNorm(z) >= gamma1-beta {
            continue // Restart
        }

        // Check hint
        r0 := lowBits(
            subVectors(w, scalarMulVector(cPoly, sk.S2)),
            alpha,
        )
        if infNorm(r0) >= gamma2-beta {
            continue // Restart
        }

        return Signature{C: c, Z: z}
    }
}
\end{lstlisting}

\subsection{Verification}

\begin{enumerate}[nosep]
  \item Compute $\mathbf{w}' = \mathbf{A}\mathbf{z} - c \cdot \mathbf{t}$.
  \item Extract high bits: $\mathbf{w}'_1$.
  \item Check: $c = H(\text{msg} \| \mathbf{w}'_1)$ and $\|\mathbf{z}\|_\infty < \gamma_1 - \beta$.
\end{enumerate}

Verification is straightforward and never requires rejection sampling.

%% ============================================================
\section{Ringtail: Threshold Lattice Signatures}

Ringtail is Lux Network's threshold signature scheme based on Ring-LWE. It allows
$t$-of-$n$ validators to collaboratively sign messages without any single validator
holding the full secret key.

\subsection{Design Challenges}

Threshold Dilithium is hard because of rejection sampling: each signer's partial
signature must independently pass the norm bound, and the combined signature must also
pass. Naive approaches have exponentially low success probability as $t$ grows.

Ringtail addresses this through:
\begin{enumerate}[nosep]
  \item \textbf{Expanded masking range}: Each signer uses a larger $\gamma_1$ to ensure
    individual rejection rarely triggers.
  \item \textbf{Aggregation-friendly bounds}: The combined norm bound accounts for the sum
    of $t$ partial signatures.
  \item \textbf{Share refresh}: Proactive secret sharing allows periodic refresh of shares
    without changing the group public key.
\end{enumerate}

\subsection{Key Distribution}

\begin{lstlisting}[language=Go,caption={Ringtail key generation via Shamir's secret sharing over $R_q$}]
func DistributeKey(n, t int) (PublicKey, []SecretShare) {
    // Generate master secret
    s := sampleSecret()

    // Shamir's secret sharing over the ring R_q
    // f(x) = s + a_1*x + ... + a_{t-1}*x^{t-1}
    coeffs := make([]RingElement, t)
    coeffs[0] = s
    for i := 1; i < t; i++ {
        coeffs[i] = sampleRing()
    }

    // Evaluate polynomial at n points
    shares := make([]SecretShare, n)
    for i := 0; i < n; i++ {
        xi := ringFromScalar(uint32(i + 1))
        shares[i] = SecretShare{
            Index: i + 1,
            Value: evaluatePoly(coeffs, xi),
        }
    }

    // Public key from master secret
    pk := computePublicKey(s)
    return pk, shares
}
\end{lstlisting}

%% ============================================================
\section{Our Go Implementation}

The entire lattice cryptography stack for Lux Network is implemented in Go with SIMD
acceleration via assembly (covered in our companion NTT paper).

\subsection{Package Structure}

\begin{lstlisting}[language=bash,caption={Package layout}]
luxfi/crypto/
    lattice/
        ring.go       # Ring element operations
        ntt.go        # NTT forward/inverse
        ntt_amd64.s   # AVX2 NTT butterfly
        poly.go       # Polynomial arithmetic
        sample.go     # Sampling (CBD, uniform, challenge)
    mlkem/
        keygen.go     # ML-KEM key generation
        encaps.go     # Encapsulation
        decaps.go     # Decapsulation
    mldsa/
        keygen.go     # ML-DSA key generation
        sign.go       # Signing with rejection sampling
        verify.go     # Verification
    ringtail/
        threshold.go  # Threshold operations
        share.go      # Secret share management
        aggregate.go  # Partial signature aggregation
\end{lstlisting}

\subsection{Performance}

\begin{table}[H]
\centering
\small
\begin{tabular}{@{}lrrl@{}}
\toprule
\textbf{Operation} & \textbf{Time} & \textbf{vs.\ Reference} & \textbf{Platform} \\
\midrule
ML-KEM-768 KeyGen & 28\,$\mu$s & 1.1x ref C & Apple M2 \\
ML-KEM-768 Encaps & 34\,$\mu$s & 1.2x ref C & Apple M2 \\
ML-KEM-768 Decaps & 38\,$\mu$s & 1.1x ref C & Apple M2 \\
ML-DSA-65 KeyGen & 52\,$\mu$s & 1.3x ref C & Apple M2 \\
ML-DSA-65 Sign & 186\,$\mu$s & 1.4x ref C & Apple M2 \\
ML-DSA-65 Verify & 56\,$\mu$s & 1.2x ref C & Apple M2 \\
NTT-256 (pure Go) & 3.1\,$\mu$s & 2.8x ref C & Apple M2 \\
NTT-256 (AVX2) & 0.9\,$\mu$s & 0.95x ref C & AMD Zen 4 \\
\bottomrule
\end{tabular}
\caption{Performance of Lux Network's Go lattice crypto implementation.}
\end{table}

%% ============================================================
\section{Security Parameters and Concrete Hardness}

\subsection{How to Choose Parameters}

Lattice parameters are chosen to resist the best known attacks. The primary attack is
the \emph{primal lattice attack} using the BKZ algorithm:

\begin{enumerate}[nosep]
  \item Construct a lattice from the LWE instance.
  \item Run BKZ with block size $\beta$ to find a short vector.
  \item The cost is approximately $2^{0.292\beta}$ classical operations (Core-SVP model).
\end{enumerate}

For 128-bit security, we need $\beta \geq 439$, which corresponds to ML-KEM-768 parameters.

\subsection{Quantum Security}

The best known quantum speedup for lattice attacks is Grover-like: $2^{0.265\beta}$
instead of $2^{0.292\beta}$. This is a modest improvement---unlike the exponential
speedup Shor provides against RSA/ECC. Lattice schemes with 128-bit classical security
provide approximately 117-bit quantum security, which is sufficient.

%% ============================================================
\section{Common Implementation Pitfalls}

\begin{enumerate}
  \item \textbf{Timing leaks in rejection sampling}: The number of loop iterations in
    ML-DSA signing depends on the secret key. Use constant-time comparison and ensure
    the loop structure does not leak via timing. Our implementation always performs the
    norm check but masks the result.

  \item \textbf{Incorrect NTT domain management}: NTT and inverse NTT must be applied in
    the right order. Adding polynomials in different domains (one in NTT form, one in
    coefficient form) is a common bug. Our type system distinguishes \texttt{Poly} from
    \texttt{NTTPoly} at the Go type level.

  \item \textbf{Missing modular reduction}: Coefficients must be reduced modulo $q$ after
    every arithmetic operation. A single missed reduction can cause subtle correctness
    failures that appear only for specific inputs.

  \item \textbf{Endianness}: NIST test vectors use little-endian byte encoding for
    polynomials. Mixing up endianness produces code that passes self-consistency tests
    but fails against reference implementations.
\end{enumerate}

%% ============================================================
\section{Conclusion}

Lattice cryptography is not magic---it is linear algebra with noise. LWE says: linear
equations are easy; noisy linear equations are hard. Ring-LWE says: we can exploit
polynomial ring structure for efficiency without losing hardness. Module-LWE says: we
can parameterize security by matrix dimension for flexibility.

The NIST standards (ML-KEM, ML-DSA) are specific instantiations with concrete parameters
chosen for practical security. Ringtail extends them to threshold operations for
decentralized signing. The NTT makes all of it fast.

This paper covered the theory. Our companion papers cover the practice: NTT with SIMD
acceleration, EVM precompiles for on-chain verification, and the formal proofs that
tie it all together.

\begin{thebibliography}{99}
\bibitem{regev2005} Regev, O. ``On Lattices, Learning with Errors, Random Linear Codes, and Cryptography.'' STOC, 2005.
\bibitem{lyubashevsky2010} Lyubashevsky, V., Peikert, C., and Regev, O. ``On Ideal Lattices and Learning with Errors over Rings.'' EUROCRYPT, 2010.
\bibitem{kyber} Avanzi, R. et al. ``CRYSTALS-Kyber: Algorithm Specifications and Supporting Documentation.'' NIST PQC Round 3, 2021.
\bibitem{dilithium} Ducas, L. et al. ``CRYSTALS-Dilithium: Algorithm Specifications and Supporting Documentation.'' NIST PQC Round 3, 2021.
\bibitem{fips203} NIST. ``Module-Lattice-Based Key-Encapsulation Mechanism Standard (ML-KEM).'' FIPS 203, 2024.
\bibitem{fips204} NIST. ``Module-Lattice-Based Digital Signature Standard (ML-DSA).'' FIPS 204, 2024.
\bibitem{shor1994} Shor, P.W. ``Algorithms for Quantum Computation: Discrete Logarithms and Factoring.'' FOCS, 1994.
\bibitem{peikert2016} Peikert, C. ``A Decade of Lattice Cryptography.'' Foundations and Trends in Theoretical Computer Science, 2016.
\end{thebibliography}

\end{document}
